\section{Single Patch Simulation Results}
\label{sec:single-patch-results}

The single inset-fed patch was simulated first because it is the radiating element used later in the
\(2\times1\) array.  The analytical dimensions in Section~\ref{sec:single-patch-analytical-design} were used as
the CST starting point, and the model was then tuned using the patch length and inset depth.  The
single-patch CST simulations were performed over the frequency range \(2.0\) to
\(3.0~\mathrm{GHz}\).  The final tuned single-patch model used PEC metallization to satisfy the CST
Learning Edition mesh-cell limitation, while the dielectric substrate loss of DiClad 880-IM was
retained through the specified loss tangent.


\subsection{Final Single-Patch Geometry Documentation}
\label{subsec:single-patch-geometry-documentation}

The final tuned geometry is documented by the top-view CST export in
Figure~\ref{fig:single-patch-top-view-dimensions} and the layered cross-sectional perspective in
Figure~\ref{fig:single-patch-cross-section-dimensions}. Because the complete CST parameter list was
longer than a single readable screen capture, the principal final dimensions are reported explicitly in
Table~\ref{tab:single-patch-final-geometry}. This combination preserves a clear CAD view while providing
an auditable dimensional record.

\begin{figure}[H]
    \centering
    \maybeincludegraphics[width=1\textwidth]{figures/5g_patch_array/single_patch/single_patch_top_view_with_dimensions.png}
    \caption[Top view of the final tuned single patch.]{Top view of the final tuned inset-fed single-patch antenna. The patch width is aligned with the CST $x$-axis, and the patch length and feed line are aligned with the $y$-axis. The final dimensions are summarized in Table~\ref{tab:single-patch-final-geometry}.}
    \label{fig:single-patch-top-view-dimensions}
\end{figure}

\begin{table}[H]
\centering
\caption{Principal final CST geometry and material parameters of the tuned single-patch antenna.}
\label{tab:single-patch-final-geometry}
\small
\setlength{\tabcolsep}{5pt}
\renewcommand{\arraystretch}{1.16}
\begin{tabularx}{0.94\textwidth}{@{}
>{\raggedright\arraybackslash}p{0.25\textwidth}
>{\centering\arraybackslash}p{0.20\textwidth}
>{\centering\arraybackslash}p{0.20\textwidth}
>{\raggedright\arraybackslash}X
@{}}
\toprule
\textbf{Quantity} & \textbf{CST parameter} & \textbf{Final value} & \textbf{Function} \\
\midrule
Target frequency & \texttt{f0} & $2.501~\mathrm{GHz}$ & Selected operating frequency \\
Relative permittivity & \texttt{epsr} & $2.20$ & DiClad 880-IM design value \\
Loss tangent & \texttt{lossTan} & $0.0009$ & Dielectric loss retained in CST \\
Substrate thickness & \texttt{h} & $1.524~\mathrm{mm}$ & Dielectric-layer thickness \\
Metal thickness & \texttt{copperT} & $0.035~\mathrm{mm}$ & Geometric thickness of each metal layer \\
Patch width & \texttt{patchW} & $47.382~\mathrm{mm}$ & Dimension along the $x$-axis \\
Patch length & \texttt{patchL} & $38.40~\mathrm{mm}$ & Final resonant dimension along the $y$-axis \\
Feed-line width & \texttt{feedW} & $4.735~\mathrm{mm}$ & Approximate $50~\Omega$ microstrip width \\
Feed-line length & \texttt{feedL} & $15.000~\mathrm{mm}$ & Feed extension from board edge to patch \\
Inset depth & \texttt{insetY} & $14.11~\mathrm{mm}$ & Final tuned feed position \\
Inset gap & \texttt{insetGap} & $1.000~\mathrm{mm}$ & Gap between feed and inset arms \\
Ground/substrate width & \texttt{gndW} & $56.526~\mathrm{mm}$ & Finite board width \\
Ground/substrate length & \texttt{gndL} & $63.801~\mathrm{mm}$ & Finite board length including feed extension \\
\bottomrule
\end{tabularx}
\end{table}

\begin{figure}[H]
    \centering
    \maybeincludegraphics[width=0.92\textwidth]{figures/5g_patch_array/single_patch/single_patch_cross_section_with_dimensions.png}
    \caption[Layered cross-sectional view of the final single-patch model.]{Layered cross-sectional perspective of the final single-patch CST model. The view shows the bottom ground metallization, DiClad 880-IM substrate, and top patch/feed metallization. The exact layer coordinates and thicknesses are listed in Table~\ref{tab:single-patch-layer-stack}.}
    \label{fig:single-patch-cross-section-dimensions}
\end{figure}

\begin{table}[H]
\centering
\caption{Layer stack used in the final single-patch CST model.}
\label{tab:single-patch-layer-stack}
\small
\renewcommand{\arraystretch}{1.16}
\begin{tabular}{@{}lccc@{}}
\toprule
\textbf{Layer} & \textbf{$z_{\min}$ (mm)} & \textbf{$z_{\max}$ (mm)} & \textbf{Thickness (mm)} \\
\midrule
Ground metallization & 0.000 & 0.035 & 0.035 \\
DiClad 880-IM substrate & 0.035 & 1.559 & 1.524 \\
Top metallization & 1.559 & 1.594 & 0.035 \\
\bottomrule
\end{tabular}
\end{table}

\subsection{Single Patch Tuning Log}
\label{subsec:single-patch-run-plan}

The patch length controls the resonant frequency, while the inset depth controls the input impedance
match.  The first simulation resonated below the target frequency, indicating that the patch was
electrically too long.  Therefore, the patch length was reduced and the inset depth was scaled to
preserve a similar feed position.  A final fine-tuning step moved the resonance very close to the
selected design frequency.

\begin{table}[H]
\centering
\caption{Single-patch CST tuning log.}
\label{tab:single-patch-tuning-log}
\footnotesize
\setlength{\tabcolsep}{4pt}
\renewcommand{\arraystretch}{1.25}

\begin{tabularx}{\textwidth}{@{}
>{\raggedright\arraybackslash}p{0.23\textwidth}
>{\raggedright\arraybackslash}p{0.30\textwidth}
>{\raggedright\arraybackslash}p{0.18\textwidth}
>{\raggedright\arraybackslash}X
@{}}
\toprule
\textbf{Run} &
\textbf{Geometry} &
\textbf{\(S_{11}\) minimum} &
\textbf{Observation} \\
\midrule

Initial analytical model &
\(L=39.657~\mathrm{mm}\),
\(y_{\mathrm{inset}}=14.520~\mathrm{mm}\),
\(g_{\mathrm{inset}}=1.000~\mathrm{mm}\) &
\(-20.39~\mathrm{dB}\) at \(2.431~\mathrm{GHz}\) &
The resonance was below the target frequency, indicating that the patch was electrically too long. \\

Retuned run &
\(L=38.55~\mathrm{mm}\),
\(y_{\mathrm{inset}}=14.11~\mathrm{mm}\),
\(g_{\mathrm{inset}}=1.000~\mathrm{mm}\) &
\(-30.70~\mathrm{dB}\) at \(2.491~\mathrm{GHz}\) &
The resonance moved close to the target frequency \(f_0=2.501~\mathrm{GHz}\). \\

Final single patch &
\(L=38.40~\mathrm{mm}\),
\(y_{\mathrm{inset}}=14.11~\mathrm{mm}\),
\(g_{\mathrm{inset}}=1.000~\mathrm{mm}\) &
\(-39.87~\mathrm{dB}\) at \(2.505~\mathrm{GHz}\) &
This final tuned result was selected for the single-patch radiation-pattern and efficiency analysis. \\

\bottomrule
\end{tabularx}
\end{table}

Figure~\ref{fig:single-patch-s11-tuning} shows the three tuning curves, and
Figure~\ref{fig:single-patch-s11-final} isolates the final tuned result.  The final resonant frequency is
\(2.50507~\mathrm{GHz}\), which is only about \(4.07~\mathrm{MHz}\) above the selected
\(2.501~\mathrm{GHz}\) design target.  This corresponds to a frequency error of approximately
\(0.16\%\), which is acceptable for the tuned CST model.

\begin{figure}[H]
    \centering
    \maybeincludegraphics[width=0.96\textwidth]{figures/5g_patch_array/single_patch/single_patch_s11_tuning_comparison.png}
    \caption{Single-patch \(S_{11}\) tuning history.  The initial analytical design resonated at approximately \(2.431~\mathrm{GHz}\).  After reducing the patch length, the resonance moved toward the target frequency, and the final tuned run resonated near \(2.505~\mathrm{GHz}\).}
    \label{fig:single-patch-s11-tuning}
\end{figure}

\begin{figure}[H]
    \centering
    \maybeincludegraphics[width=0.96\textwidth]{figures/5g_patch_array/single_patch/single_patch_s11_tuned_final.png}
    \caption{Final tuned reflection coefficient of the single inset-fed patch antenna.  The final model resonated at approximately \(2.50507~\mathrm{GHz}\) with \(S_{11,\min}\approx -39.87~\mathrm{dB}\), very close to the selected design frequency of \(2.501~\mathrm{GHz}\).}
    \label{fig:single-patch-s11-final}
\end{figure}

\subsection{Single Patch Radiation Patterns and Efficiency}

The final single-patch farfield result is shown in Figure~\ref{fig:single-patch-3d-farfield}.  CST
reported a maximum directivity of approximately \(7.226~\mathrm{dBi}\) at
\(2.501~\mathrm{GHz}\).  The radiation efficiency and total efficiency were reported as
\(-0.1634~\mathrm{dB}\) and \(-0.1700~\mathrm{dB}\), respectively, corresponding to approximately
\(96.3\%\) and \(96.2\%\).  Because PEC metallization was used for the tuned run, the radiation efficiency includes
the modeled dielectric loss but excludes conductor loss and input-mismatch loss. The total
efficiency additionally includes input-mismatch loss. Neither value includes finite-copper ohmic
loss.

\begin{figure}[H]
    \centering
    \maybeincludegraphics[width=0.88\textwidth]{figures/5g_patch_array/single_patch/single_patch_3d_farfield_tuned.png}
    \caption{Three-dimensional CST farfield pattern of the final tuned single patch at \(2.501~\mathrm{GHz}\).  The reported maximum directivity is approximately \(7.226~\mathrm{dBi}\), with radiation and total efficiencies of approximately \(96.3\%\) and \(96.2\%\), respectively.}
    \label{fig:single-patch-3d-farfield}
\end{figure}

Figures~\ref{fig:single-patch-phi0} and~\ref{fig:single-patch-phi90} show the two principal 2D
farfield cuts.  The \(\phi=0^\circ\) cut gives a main-lobe magnitude of approximately
\(7.23~\mathrm{dBi}\), a broadside main-lobe direction of \(0^\circ\), and a 3-dB angular width of
approximately \(86.8^\circ\).  These values are consistent with the expected broadside radiation of a
rectangular microstrip patch.

\begin{figure}[H]
    \centering
    \maybeincludegraphics[width=0.86\textwidth]{figures/5g_patch_array/single_patch/single_patch_2d_pattern_phi0_tuned.png}
    \caption{Single-patch farfield directivity cut for \(\phi=0^\circ\).  The main lobe is broadside with a magnitude of about \(7.23~\mathrm{dBi}\) and an angular width of about \(86.8^\circ\).}
    \label{fig:single-patch-phi0}
\end{figure}

\begin{figure}[H]
    \centering
    \maybeincludegraphics[width=0.86\textwidth]{figures/5g_patch_array/single_patch/single_patch_2d_pattern_phi90_tuned.png}
    \caption{Single-patch farfield directivity cut for \(\phi=90^\circ\).  This cut provides the orthogonal principal-plane response of the tuned patch element.}
    \label{fig:single-patch-phi90}
\end{figure}

\begin{figure}[H]
    \centering
    \maybeincludegraphics[width=0.92\textwidth]{figures/5g_patch_array/single_patch/single_patch_efficiencies_tuned.png}
    \caption[Efficiency of the final tuned single-patch PEC model.]{CST efficiency result for the final tuned single-patch PEC model at \(2.501~\mathrm{GHz}\). The simulated radiation efficiency is approximately \(-0.1634~\mathrm{dB}\) \(\approx96.3\%\), and the simulated total efficiency is approximately \(-0.1700~\mathrm{dB}\) \(\approx96.2\%\).}
    \label{fig:single-patch-efficiency}
\end{figure}

\begin{table}[H]
\centering
\caption{Single patch final simulation summary.}
\label{tab:single-patch-final-summary}
\begin{tabular}{@{}ll@{}}
\toprule
\textbf{Metric} & \textbf{Final CST Value} \\
\midrule
Target frequency & \(2.501~\mathrm{GHz}\) \\
Final resonant frequency & \(2.50507~\mathrm{GHz}\) \\
Frequency error & \(0.16\%\) \\
Minimum \(S_{11}\) & \(-39.87~\mathrm{dB}\) \\
Radiation efficiency & \(-0.1634~\mathrm{dB}\) \(\approx 96.3\%\) \\
Total efficiency & \(-0.1700~\mathrm{dB}\) \(\approx 96.2\%\) \\
Directivity & \(7.226~\mathrm{dBi}\) \\
Main-beam direction, \(\phi=0^\circ\) cut & \(0^\circ\) \\
3-dB angular width, \(\phi=0^\circ\) cut & \(86.8^\circ\) \\
Final tuned \texttt{patchL} & \(38.40~\mathrm{mm}\) \\
Final tuned \texttt{insetY} & \(14.11~\mathrm{mm}\) \\
\bottomrule
\end{tabular}
\end{table}

\subsection{Finite-Conductivity Copper Validation}
\label{subsec:single-patch-copper-validation}

A supplementary conductor-loss validation was performed on the final single-patch geometry after the
PEC baseline had been completed. The geometry, substrate, port, boundaries, solver type, frequency
range, and tuned dimensions were held fixed. The only physical change was replacement of the PEC
metallization by the CST \emph{Lossy metal} model with
\[
\sigma_{\mathrm{Cu}}=5.87\times10^7~\mathrm{S/m}.
\]
This produced a controlled comparison of the PEC approximation against a finite-conductivity copper
model.

An initial lossy-metal run with adaptive tetrahedral mesh refinement was attempted. The adaptive
sequence reached the fourth refinement pass, but the mesh then exceeded the
\(20{,}000\)-cell limit of the CST Learning Edition before a converged solution was obtained. The
incomplete adaptive result was discarded. The simulation was then repeated using the same geometry and fixed initial-mesh settings, with mesh adaptation disabled. This fixed-mesh run completed successfully after eight
frequency samples and produced valid \(S\)-parameters, impedance matrices, VSWR data, efficiencies,
and farfield results. Consequently, the study is treated as a fixed-mesh conductor-model sensitivity
check rather than a formal mesh-convergence study.

Figure~\ref{fig:single-patch-copper-s11} shows that finite copper conductivity had very little effect on
the tuned resonance and impedance match. The lossy-metal model resonated near
\(2.504~\mathrm{GHz}\) with
\(S_{11,\min}\approx-39.72~\mathrm{dB}\), compared with
\(2.50507~\mathrm{GHz}\) and \(-39.87~\mathrm{dB}\) for the PEC baseline.

\begin{figure}[H]
    \centering
    \maybeincludegraphics[width=0.94\textwidth]{figures/5g_patch_array/single_patch/single_patch_copper_lossymetal_s11.png}
    \caption[Finite-conductivity single-patch reflection coefficient.]{Reflection coefficient of the fixed-mesh finite-conductivity copper model. The resonance occurs near \(2.504~\mathrm{GHz}\) with \(S_{11,\min}\approx-39.72~\mathrm{dB}\), closely matching the PEC baseline.}
    \label{fig:single-patch-copper-s11}
\end{figure}

The finite-conductivity model produced a radiation efficiency of
\(-0.340329~\mathrm{dB}\), corresponding to approximately \(92.46\%\), and a total efficiency of
\(-0.344040~\mathrm{dB}\), corresponding to approximately \(92.38\%\). The small difference between
the two efficiencies is consistent with the very strong impedance match near the operating frequency.
The efficiency reduction relative to PEC quantifies the conductor-loss contribution that was omitted
from the baseline model.

\begin{figure}[H]
    \centering
    \begin{subfigure}[t]{0.49\textwidth}
        \centering
        \maybeincludegraphics[width=\textwidth]{figures/5g_patch_array/single_patch/single_patch_copper_lossymetal_radiation_efficiency.png}
        \caption{Radiation efficiency.}
        \label{fig:single-patch-copper-rad-eff}
    \end{subfigure}
    \hfill
    \begin{subfigure}[t]{0.49\textwidth}
        \centering
        \maybeincludegraphics[width=\textwidth]{figures/5g_patch_array/single_patch/single_patch_copper_lossymetal_total_efficiency.png}
        \caption{Total efficiency.}
        \label{fig:single-patch-copper-tot-eff}
    \end{subfigure}
    \caption[Finite-conductivity single-patch efficiencies.]{Efficiencies of the fixed-mesh finite-conductivity copper model at \(2.501~\mathrm{GHz}\). The radiation efficiency is approximately \(92.46\%\), and the total efficiency is approximately \(92.38\%\).}
    \label{fig:single-patch-copper-efficiencies}
\end{figure}

The peak directivity remained approximately \(7.23~\mathrm{dBi}\), essentially unchanged from the
PEC value of \(7.226~\mathrm{dBi}\). The \(\phi=90^\circ\) principal-plane cut retained a broadside
main lobe near \(2^\circ\), a 3-dB beamwidth of approximately \(80.8^\circ\), and a sidelobe level of
approximately \(-21.2~\mathrm{dB}\). The farfield comparison therefore indicates that conductor loss
reduced efficiency without materially changing the normalized radiation-pattern shape or directivity.

\begin{figure}[H]
    \centering
    \begin{subfigure}[t]{0.47\textwidth}
        \centering
        \maybeincludegraphics[width=\textwidth]{figures/5g_patch_array/single_patch/single_patch_copper_lossymetal_3d_farfield.png}
        \caption{Three-dimensional directivity pattern.}
        \label{fig:single-patch-copper-3d}
    \end{subfigure}
    \hfill
    \begin{subfigure}[t]{0.51\textwidth}
        \centering
        \maybeincludegraphics[width=\textwidth]{figures/5g_patch_array/single_patch/single_patch_copper_lossymetal_phi90_pattern.png}
        \caption{Principal-plane cut for \(\phi=90^\circ\).}
        \label{fig:single-patch-copper-phi90}
    \end{subfigure}
    \caption[Finite-conductivity single-patch radiation patterns.]{Radiation patterns of the fixed-mesh finite-conductivity copper model at \(2.501~\mathrm{GHz}\). The peak directivity remains approximately \(7.23~\mathrm{dBi}\), and the principal-plane pattern remains consistent with the PEC result.}
    \label{fig:single-patch-copper-patterns}
\end{figure}

\begin{table}[H]
\centering
\caption{PEC and finite-conductivity copper comparison for the final single-patch geometry.}
\label{tab:single-patch-conductor-comparison}
\small
\setlength{\tabcolsep}{4pt}
\renewcommand{\arraystretch}{1.18}
\begin{tabularx}{\textwidth}{@{}
>{\raggedright\arraybackslash}X
>{\centering\arraybackslash}p{0.22\textwidth}
>{\centering\arraybackslash}p{0.25\textwidth}
>{\raggedright\arraybackslash}p{0.28\textwidth}
@{}}
\toprule
\textbf{Metric} & \textbf{PEC baseline} & \textbf{Copper lossy metal} & \textbf{Engineering interpretation} \\
\midrule
Resonant frequency & \(2.50507~\mathrm{GHz}\) & \(2.504~\mathrm{GHz}\) & Change of approximately \(-1.07~\mathrm{MHz}\) \\
Minimum \(S_{11}\) & \(-39.87~\mathrm{dB}\) & \(-39.72~\mathrm{dB}\) & Matching essentially unchanged \\
Peak directivity & \(7.226~\mathrm{dBi}\) & \(\approx7.23~\mathrm{dBi}\) & Directivity essentially unchanged \\
Radiation efficiency & \(96.3\%\) & \(92.46\%\) & Reduction of approximately \(3.84\) percentage points \\
Total efficiency & \(96.2\%\) & \(92.38\%\) & Reduction of approximately \(3.78\) percentage points \\
\(\phi=90^\circ\) 3-dB beamwidth & \(80.8^\circ\) & \(80.8^\circ\) & Normalized pattern remained stable \\
\bottomrule
\end{tabularx}
\end{table}

The comparison confirms that the PEC approximation was accurate for resonant frequency, impedance
matching, directivity, and radiation-pattern shape in this project. Its principal limitation was an
overestimate of radiation and total efficiency because finite copper ohmic loss was excluded.
